\scanpage{38}\mspage{37}

satisfaisant

a) $f_{\unsure{\mathrm{str}|\mathrm{str.}}}$ croissante pour $\leq$ et $\ll$

b) $\deleted{\unsure{f_{\mathrm{str}}}}$ \deleted{\unsure{pour tout|pour toute}} $F\subset M$ provient d'une figure, $f_{\mathrm{str}}(F)$ est une figure

\margin{\tikz[baseline=(ms38NB.base)]{\node[inner sep=.5pt] (ms38NB) {\uline{NB}}; \draw ([xshift=2pt]ms38NB.north east) -- (ms38NB.north west) -- (ms38NB.south west) -- ([xshift=1pt]ms38NB.south east);} Pour que $f$ soit inj. il f et s que $f_{\mathrm{str}}$ le soit. Pour que ce soit un iso. avec un sous-atelier \deleted{strict}, il f. et s. que $f$ soit un plongement pour $\leq$ et $\ll$, et que l'image soit $\leq$-fermée.$^{*}$}
% slanted
\margin{Mais le sous-atelier n'est pas strict pour autant -- dire qu'il l'est \unsure{signifie} que toute partie de $f(M)\subset M'$ qui est une figure pour $\unsure{\mathcal F|\mathscr F}\mathstrut'$, en est déjà une pour $\mathcal F$}

On appelle \uline{sous-atelier} d'un atelier, une \deleted{\unsure{de}} partie $\mathcal F'$ de $\mathcal F$, telle que, munie des relations d'ordre $\underset{\mathcal F'}{\leq},\underset{\mathcal F}{\ll}$ induites par celles de $\mathcal F$, ce soit un atelier, et que $\mathcal F'\hookrightarrow\mathcal F$ soit un morphisme d'ateliers, i.e. que

$1^{\circ}$) Si $F_i\in\mathcal F'$ sont majorés \uline{dans} $\mathcal F'$, \inserted{pour $\leq$} alors
\[
\mathop{\mathrm{Sup}}\limits_{i\,(\mathcal F)}F_i\in\mathcal F'\qquad\text{(stabilité par Sup majorés)}
\]
$2^{\circ}$) Sous \deleted{\unsure{la}} les $=$ hypothèses, $\bigcap F_i\in\mathcal F'$

% slanted
\margin{Mais on ne suppose pas que, si $F,F'\in\mathcal F'$ et $F\vee F'$ existe dans $\mathcal F$ (i.e. $\{F,F'\}$ majorés dans $\mathcal F$) on ait $F\vee F'\in\mathcal F'$ i.e. $\{F,F'\}$ majoré dans $\mathcal F'$. Donc on ne suppose \deleted{pas} $\mathcal F'$ stable par sommes (quand elles \unsure{ne le sont|ne sont} majorées dans $\mathcal F'$) Le sous-atelier est dit (\uline{\unsure{spacieux}} ou \uline{strict}) s'il est stable pour Sup dans $\mathcal F$ (\unsure{pas néc.|pour néc.} majorés)}

\tikz[remember picture,overlay]\coordinate (msQuestionTop);
\unsure{Je} me demande s'il ne serait pas plus raisonnable, au lieu de $1^{\circ}$, et $2^{\circ}$, de poser

$\text{\textcircled{\scriptsize $*$}}$\quad \uline{Si} $F'\in\mathcal F'$, \deleted{et $F\in\mathcal F$ est tel que $F$} \inserted{\uline{alors \unsure{tte|toute} sous figure}} $F'$ \uline{est dans} $\mathcal F'$.

Il faut vérifier \inserted{seulement} que cette condition suffit à impliquer que $(\mathcal F',\underset{\mathcal F'}{\leq},\underset{\mathcal F'}{\ll})$ est bien un atelier\,\ldots\tikz[remember picture,overlay]\coordinate (msQuestionEnd);
\figure{\tikz[remember picture,overlay]{
  \draw ([yshift=6pt]msQuestionTop) -- ([xshift=-4pt,yshift=6pt]msQuestionTop)
    -- ([xshift=-4pt,yshift=-2pt]msQuestionTop |- msQuestionEnd) -- ([yshift=-2pt]msQuestionTop |- msQuestionEnd);
}}

Les sous-ateliers \inserted{stricts} correspondent aux sous ens. \inserted{$M'$} de $M$ fermés pour $\leq$, sans plus,
% slanted
\margin{En tous les cas, si $\mathcal F'$ est un sous-atelier de $\mathcal F$, $\mathcal F'\unsure{\rule[.5ex]{1.4em}{.55pt}|\longrightarrow}\mathcal F$ est \deleted{\unsure{un}} un morphisme strict et $M'$ s'identifie à l'ens des fig. él. dans $\mathcal F$. (et $M'=\mathcal F'\cap M$?)}

Donc, $\mathcal F$ est dit ensembliste, si l'application $\mathrm{mulomb}\colon\mathcal F\unsure{\rule[.5ex]{1.4em}{.55pt}|\longrightarrow}\mathrm{Fig}(L)$ le ``réalise'' comme sous-\unsure{atelier} de $\mathrm{Fig}\,L$ -- en fait, on aimerait que ce soit un sous-atelier \uline{strict}$^{*}$,
% slanted
\margin{$^{*}$ C'est donc une \unsure{condition} plus forte encore que dans la définition de la page précédente, où d) implique seulement la stabilité pour $F\vee G$, et pas pour les Sup quelconques.}
et on voudra \deleted{que} en particulier que $M$ se réalise comme un ens de figures élémentaires (ou multistrates ensemblistes) \uline{dans $L$}. Ainsi, dans la définition de la page précédente, il faudra \unsure{rajouter|se contenter|se pencher}
